% =====================================================================================================
% PART 4 (frames 15-18): more colours -- the soundness criterion, why it holds, the family for
% k = 3 + L colours, universal vertices and perfect completeness.
% Sources in the paper: Section 1.2 (paragraphs "A soundness criterion", "Universal vertices"),
% Section 4 (Theorem 4.1, proof idea, Lemma 4.2, Proposition 4.3), Section 5 (the family, Lemma 5.1,
% proof of Theorem 1.2, Figure 4 (c),(d)), Section 6 (Lemma 6.1, proof of Theorem 1.3, Remark 6.2),
% Table 1.
% Figures: figs/p4-decoding.tex, figs/p4-gadgets-AN.tex, figs/p4-universal.tex
% =====================================================================================================

% local helpers of part 4 (names are prefixed with pIV to avoid clashes with other parts)
\providecommand{\pIVto}{\ensuremath{\longrightarrow}}
\providecommand{\pIVbox}[1]{\tikz[baseline=(b.base)]\node[draw=acc,line width=0.8pt,fill=acc!5,
  rounded corners=3pt,inner xsep=6pt,inner ysep=5.5pt] (b) {#1};}

% ---- Frame 15: More colours: a general soundness criterion
\begin{frame}{More colours: a general soundness criterion}
\vspace{-0.5mm}
$k\ge4$:\ satisfying colourings use only $3$ of the $k$ colours\ \ $\Rightarrow$\ \
\hl{global vertices} needed again
\par\vspace{2.4mm}
\begin{block}{Theorem 4.1\ \ (soundness criterion)}
Gadgets $I_0,I_1,I_2$ for \MC{k}, one per type
{\small\color{soft}(each symmetric in the 2nd and 3rd variable)}.\\[0.8mm]
Then, for every $\delta>0$, \MC{k} is NP-hard to approximate within $\rho+\delta$, where
\par\vspace{1.4mm}
\centerline{%
\begin{tikzpicture}[baseline=(f.base)]
\node[inner sep=1pt] (f) {\large$\displaystyle\rho\;=\;
  \frac{{\color{acc}\max\limits_{x,\,y}}\ \sum\limits_{t}V_t(x,y)}{\sum\limits_{t}C_t}$};
\node[anchor=west,font=\small,text=acc,inner sep=1.5pt,align=left] (n) at ([xshift=13mm,yshift=4.4mm]f.east)
  {one colour pattern $(x,y)$,\\[-0.4mm]all three gadgets};
\draw[-{Stealth[length=4pt,width=3.4pt]},line width=0.5pt,acc] (n.west) -- ([xshift=1.5mm,yshift=4.4mm]f.east);
\end{tikzpicture}}
\vspace{-1.5mm}
\end{block}
\vspace{1.4mm}
{\small
\begin{tabular}{@{\hspace{2mm}}r@{\quad}l@{}}
$V_t(x,y)$ & best cut of the type-$t$ gadget when its variables have colours $x$, the globals colours $y$\\[0.8mm]
$C_t$ & its least value on satisfying colourings\ \ {\color{soft}(globals coloured canonically)}
\end{tabular}}
\par\vspace{3.4mm}
\centerline{A cheating colouring cannot do better than to use}
\vspace{0.6mm}
\centerline{\large\hl{one colour pattern on all equations}}
\end{frame}

% ---- Frame 16: Why the criterion holds
\begin{frame}{Why the criterion holds}
\vspace{-1mm}
\begin{columns}[c,onlytextwidth]
\begin{column}{0.455\textwidth}
% Frame 16: decoding a colouring into an assignment (schematic; cf. the paper's Figure 2 with its labels
% "decoded by a" / "decoded by b", and the proof idea of Theorem 4.1).
% Top row: the variable vertices of U (two) and of W (four) with their colours, drawn as fills
% (k = 4 colours: white, grey, accent blue, dark grey).  Every U-vertex is sent through the map a, every
% W-vertex through the map b (two arbitrary maps [k] -> Z_3, shown as tables); the bottom row is the
% decoded assignment  u -> a(colour of u),  w -> b(colour of w).
% The two maps are different: e.g. blue is decoded to 2 on U but to 1 on W.
\begingroup
\newcommand{\pIVent}[4]{% #1,#2 = position of the swatch, #3 = colour style, #4 = value
  \node[sw,#3] at (#1,#2) {};
  \node[font=\scriptsize,inner sep=0pt] at ($(#1,#2)+(0.32,0)$) {$\mapsto$};
  \node[font=\small,inner sep=0pt] at ($(#1,#2)+(0.60,0)$) {$#4$};}
\begin{tikzpicture}[x=1cm,y=1cm,
  v/.style={circle,draw,line width=0.6pt,inner sep=0pt,minimum size=12pt},
  sw/.style={circle,draw,line width=0.45pt,inner sep=0pt,minimum size=7.5pt},
  c0/.style={fill=white}, c1/.style={fill=black!35}, c2/.style={fill=acc}, c3/.style={fill=black!78},
  grp/.style={font=\normalsize,text=black!70,inner sep=0pt,anchor=base},
  valn/.style={font=\large,inner sep=0pt,anchor=base},
  mapbox/.style={draw=acc!55,fill=acc!5,rounded corners=3pt,line width=0.6pt},
  mapname/.style={font=\Large,text=acc,inner sep=0pt},
  ar/.style={-{Stealth[length=4pt,width=3.6pt]},line width=0.6pt,black!60}]
% ---------------------------------------------------------------- the two maps (boxes)
\draw[mapbox] (0.00,1.62) rectangle (2.66,2.86);
\draw[mapbox] (2.98,1.62) rectangle (6.30,2.86);
\node[mapname] at (0.36,2.24) {$a$};
\pIVent{0.88}{2.50}{c0}{1}
\pIVent{1.88}{2.50}{c1}{0}
\pIVent{0.88}{1.98}{c2}{2}
\pIVent{1.88}{1.98}{c3}{1}
\node[mapname] at (3.52,2.24) {$b$};
\pIVent{4.26}{2.50}{c0}{0}
\pIVent{5.26}{2.50}{c1}{1}
\pIVent{4.26}{1.98}{c2}{1}
\pIVent{5.26}{1.98}{c3}{2}
% ---------------------------------------------------------------- vertices (colours), arrows, decoded values
\foreach \x/\c/\val in {0.90/c2/2, 1.76/c1/0, 3.35/c0/0, 4.21/c2/1, 5.07/c1/1, 5.93/c3/2}{
  \node[v,\c] (n) at (\x,3.50) {};
  \draw[ar] (\x,3.24) -- (\x,2.89);
  \draw[ar] (\x,1.62) -- (\x,1.22);
  \node[valn] at (\x,0.86) {$\val$};}
% ---------------------------------------------------------------- group labels
\node[grp] at (1.33,4.02) {$U$};
\node[grp] at (4.64,4.02) {$W$};
\node[font=\small,inner sep=0pt,anchor=base] at (3.405,4.62) {a \hl{colouring} of the variables\ \ {\color{soft}($k=4$)}};
% ---------------------------------------------------------------- the result
\draw[decorate,decoration={brace,mirror,amplitude=4pt},black!55,line width=0.5pt] (0.60,0.66) -- (6.21,0.66);
\node[font=\small,inner sep=0pt,anchor=base] at (3.405,0.10) {an \hl{assignment} of the equations};
\end{tikzpicture}%
\endgroup

\end{column}
\begin{column}{0.525\textwidth}
\begin{itemize}\setlength{\itemsep}{2mm}
\item Colouring\ \ \pIVto\ \ distribution $P(t,x)$ of (type, colour pattern)
\item Any maps $a,b\colon[k]\to\Z_3$ \hl{decode} it into an assignment
\item \bad{\textbf{NO case:}} it satisfies $\le\frac13+\eps$\\
  $\Rightarrow$\ every \hl{test} $(a,b)$ sees $t$ as nearly uniform and independent of $x$
\item \hl{Spanning lemma}: these tests detect every dependence that can affect the cut\\[0.3mm]
  {\small\color{soft}short Fourier argument}\hfill\paperref{Lemma 4.2}
\end{itemize}
\end{column}
\end{columns}
\vspace{2mm}
{\color{acc!35}\rule{\linewidth}{0.5pt}}\par
\vspace{1.2mm}
{\small
\hl{Needs role separation:}\ \ $U$ and $W$ are decoded by independent maps $a$ and $b$\par
\vspace{1mm}
\hl{Tight:}\ \ no black-box use of H\aa stad's theorem can improve the numerator of $\rho$\hfill\paperref{Proposition 4.3}\par}
\end{frame}

% ---- Frame 17: A uniform family for k = 3 + L colours
\begin{frame}{A uniform family for $k=3+L$ colours}
\vspace{-0.8mm}
% lead line (Section 1.2, "A soundness criterion"; Remark "How the family was found")
Criterion\ $\Rightarrow$\ gadget design is \hl{one linear program}%
\hfill{\small\color{soft}solved for $k=4,5$, simplified by hand}
\par\vspace{3.2mm}
\begin{columns}[c,onlytextwidth]
\begin{column}{0.53\textwidth}
% Frame 17: the gadgets A_L (top) and N_L (bottom), drawn for L = 2 = panels (c) and (d) of the paper's
% Figure 4 (../paper/figures/gadgets.tex), with EXACTLY the same edges and weights.  The coordinates of the
% two panels are those of the paper's figure; each panel is only scaled as a whole (x=..,y=.. of its
% scope), the vertex labels are larger, and the panel captions are replaced by slide text at the left
% (names, types, a legend of the two edge weights; the two legend strokes are not edges of a gadget).
% Every edge is drawn explicitly, one \draw per edge (the edge list is extracted from these lines and
% compared with the definitions in Section 5 of the paper).
% Edge styles: e1, e1p, u1 = weight 1;  e3, u3 = weight 3;  u1, u3 = edges at universal vertices z_i.
\begin{tikzpicture}[x=1cm,y=1cm,
  prim/.style={circle,draw,line width=0.5pt,fill=black!15,inner sep=0pt,minimum size=13pt,font=\footnotesize},
  aux/.style={circle,draw,line width=0.5pt,fill=white,inner sep=0pt,minimum size=12.5pt,font=\scriptsize},
  univ/.style={rectangle,draw,line width=0.5pt,fill=black!5,inner sep=0pt,minimum size=12.5pt,font=\scriptsize},
  e1/.style={line width=0.5pt},
  e1p/.style={line width=0.5pt,black!60},
  e3/.style={line width=1.6pt},
  u1/.style={line width=0.5pt,acc},
  u3/.style={line width=1.6pt,acc},
  gname/.style={font=\Large,text=acc,inner sep=0pt,anchor=base west},
  gsub/.style={font=\small,inner sep=0pt,anchor=base west},
  gnote/.style={font=\footnotesize,inner sep=0pt,anchor=base west,text=soft}]
% ================================================================== (c) A_L, L=2
\begin{scope}[shift={(4.72,0)},x=0.69cm,y=0.69cm]
  \node[aux] (aC1) at (90:1.7) {$a_1$};
  \node[aux] (aC4) at (30:1.7) {$a_4$};
  \node[aux] (aC5) at (-30:1.7) {$a_5$};
  \node[aux] (aC2) at (-90:1.7) {$a_2$};
  \node[aux] (aC3) at (-150:1.7) {$a_3$};
  \node[aux] (aC6) at (150:1.7) {$a_6$};
  \node[prim] (pC1) at (90:0.765) {$p_1$};
  \node[prim] (pC2) at (-30:0.765) {$p_2$};
  \node[prim] (pC3) at (210:0.765) {$p_3$};
  \node[univ] (zC1) at (-3.1,0) {$z_1$};
  \node[univ] (zC2) at (3.1,0) {$z_2$};
  % the gadget E, unchanged
  \draw[e1] (aC1) -- (aC4);
  \draw[e1] (aC4) -- (aC5);
  \draw[e1] (aC5) -- (aC2);
  \draw[e1] (aC2) -- (aC3);
  \draw[e1] (aC3) -- (aC6);
  \draw[e1] (aC6) -- (aC1);
  \draw[e1p] (pC1) -- (aC1);
  \draw[e1p] (pC1) -- (aC2);
  \draw[e1p] (pC1) -- (aC3);
  \draw[e1p] (pC1) -- (aC4);
  \draw[e1p] (pC2) -- (aC3);
  \draw[e1p] (pC2) -- (aC4);
  \draw[e1p] (pC2) -- (aC5);
  \draw[e1p] (pC2) -- (aC6);
  \draw[e1p] (pC3) -- (aC1);
  \draw[e1p] (pC3) -- (aC2);
  \draw[e1p] (pC3) -- (aC5);
  \draw[e1p] (pC3) -- (aC6);
  % z_1, z_2 -- every a_j (weight 1): straight to the nearest two a_j, curved around the cycle to the others
  \draw[u1] (zC1) -- (aC6);
  \draw[u1] (zC1) -- (aC3);
  \draw[u1] (zC2) -- (aC4);
  \draw[u1] (zC2) -- (aC5);
  % (curves are kept out of the bounding box, which would otherwise include their control points;
  %  the \path below adds their true extent)
  \begin{pgfinterruptboundingbox}
  \draw[u1] (zC1) .. controls (-2.43,0.74) and (-1.09,1.85) .. (aC1);
  \draw[u1] (zC1) .. controls (-2.43,-0.74) and (-1.09,-1.85) .. (aC2);
  \draw[u1] (zC1) .. controls (-2.47,2.52) and (0.37,3.21) .. (aC4);
  \draw[u1] (zC1) .. controls (-2.47,-2.52) and (0.37,-3.21) .. (aC5);
  \draw[u1] (zC2) .. controls (2.43,0.74) and (1.09,1.85) .. (aC1);
  \draw[u1] (zC2) .. controls (2.43,-0.74) and (1.09,-1.85) .. (aC2);
  \draw[u1] (zC2) .. controls (2.47,2.52) and (-0.37,3.21) .. (aC6);
  \draw[u1] (zC2) .. controls (2.47,-2.52) and (-0.37,-3.21) .. (aC3);
  % z_1 -- z_2 (weight 3)
  \draw[u3] (zC1) .. controls (-3.71,3.47) and (3.71,3.47) .. (zC2);
  \end{pgfinterruptboundingbox}
  \path (-3.16,-2.36) rectangle (3.16,2.68);
\end{scope}
% slide text of panel (c)
\node[gname] at (0,1.20) {$A_L$};
\node[gsub]  at (0,0.66) {type $0$};
\node[gsub]  at (0,0.18) {$=E+{}$the $z_i$};
\draw[e3] (0.02,-0.52) -- (0.50,-0.52);
\node[gnote,anchor=west] at (0.64,-0.52) {weight $3$};
\draw[e1] (0.02,-0.92) -- (0.50,-0.92);
\node[gnote,anchor=west] at (0.64,-0.92) {weight $1$};
% ================================================================== (d) N_L, L=2
\begin{scope}[shift={(4.72,-2.83)},x=0.74cm,y=0.74cm]
  \node[prim] (pD1) at (90:1.2) {$p_1$};
  \node[prim] (pD2) at (-30:1.2) {$p_2$};
  \node[prim] (pD3) at (210:1.2) {$p_3$};
  \node[aux] (qD) at (0,0) {$q$};
  % z_i raised above the height of q, so that z_1--q--z_2 is not a straight line (no edge z_1z_2)
  \node[univ] (zD1) at (-2.6,0.6) {$z_1$};
  \node[univ] (zD2) at (2.6,0.6) {$z_2$};
  % the gadget N, unchanged
  \draw[e1] (pD1) -- (pD2);
  \draw[e1] (pD2) -- (pD3);
  \draw[e1] (pD3) -- (pD1);
  \draw[e3] (qD) -- (pD1);
  \draw[e3] (qD) -- (pD2);
  \draw[e3] (qD) -- (pD3);
  % z_i -- q and z_i -- every primary (weight 3); no edge z_1 z_2
  \draw[u3] (zD1) -- (qD);
  \draw[u3] (zD1) -- (pD1);
  \draw[u3] (zD1) -- (pD3);
  \draw[u3] (zD2) -- (qD);
  \draw[u3] (zD2) -- (pD1);
  \draw[u3] (zD2) -- (pD2);
  \begin{pgfinterruptboundingbox}
  \draw[u3] (zD1) .. controls (-2.4,-1.6) and (-0.15,-1.75) .. (pD2);
  \draw[u3] (zD2) .. controls (2.4,-1.6) and (0.15,-1.75) .. (pD3);
  \end{pgfinterruptboundingbox}
  \path (-2.63,-1.38) rectangle (2.63,0);
\end{scope}
% slide text of panel (d)
\node[gname] at (0,-2.06) {$N_L$};
\node[gsub]  at (0,-2.60) {types $1$ and $2$};
\node[gsub]  at (0,-3.08) {$=N+{}$the $z_i$};
\node[gnote] at (0,-3.62) {drawn for $L=2$};
\end{tikzpicture}

\end{column}
\begin{column}{0.45\textwidth}
\raggedright
Add $L$ \hl{universal vertices} $z_1,\dots,z_L$\\
(squares):\ \ globals for the extra colours\par
\vspace{4mm}
The family $(2A_L,\,N_L,\,N_L)$ gives\par
\vspace{1.8mm}
\pIVbox{$\begin{aligned}
\rho\;&=\;1-2/\Gamma_L\\[0.4mm]
\Gamma_L\;&=\;29+18L+3\tbinom L2
\end{aligned}$}\hfill
{\color{soft}\footnotesize\begin{tabular}{@{}r@{}}(Lemma 5.1,\\[-0.4mm]by hand)\end{tabular}}\par
\vspace{4mm}
{\renewcommand{\arraystretch}{1.15}%
\begin{tabular}{@{}l@{\hspace{4mm}}c@{\hspace{4mm}}c@{\hspace{4mm}}c@{\hspace{4mm}}c@{\hspace{4mm}}c@{}}
\toprule
$k$ & $3$ & \hl{$4$} & \hl{$5$} & $6$ & $7$\\
\midrule
factor\rule[-3.6mm]{0pt}{9.4mm} & $\dfrac{27}{29}$ & \hl{$\dfrac{45}{47}$} & \hl{$\dfrac{33}{34}$} & $\dfrac{45}{46}$
  & $\dfrac{117}{119}$\\
\bottomrule
\end{tabular}}\par
\end{column}
\end{columns}
\end{frame}

% ---- Frame 18: Universal vertices and perfect completeness
\begin{frame}{Universal vertices and perfect completeness}
\vspace{-1mm}
\begin{columns}[c,onlytextwidth]
\begin{column}{0.42\textwidth}
\centering
A classical reduction\ \ \src{KKLP97, GS13}\par
\vspace{2.2mm}
% Frame 18: the universal-vertex reduction of [KKLP97, GS13] (Lemma 6.1 of the paper), schematic, L = 2.
% A base graph G (blob; a hard instance for 3 colours) and two new vertices z_1, z_2 (squares).
% Every z_i is joined to every vertex v of G (weight d_v/3, blue fan), and z_1 and z_2 are joined to each
% other.  The small graph inside the blob is only an illustration (7 vertices, 10 edges, 3-colourable).
\begin{tikzpicture}[x=1cm,y=1cm,
  gv/.style={circle,draw,line width=0.5pt,fill=white,inner sep=0pt,minimum size=7.5pt},
  univ/.style={rectangle,draw,line width=0.6pt,fill=black!5,inner sep=0pt,minimum size=15pt,font=\small},
  ge/.style={line width=0.55pt},
  fan/.style={line width=0.45pt,acc!70},
  zz/.style={line width=1.6pt,acc},
  note/.style={font=\small,inner sep=0pt,text=acc,align=left,anchor=west},
  ptr/.style={-{Stealth[length=3.5pt,width=3pt]},line width=0.4pt,acc}]
% ---------------------------------------------------------------- the base graph G
\draw[draw=black!45,fill=black!4,line width=0.6pt] (0,0) ellipse [x radius=2.55, y radius=1.0];
\coordinate (v1) at (-1.92,0.02);
\coordinate (v2) at (-1.15,0.45);
\coordinate (v3) at (-1.00,-0.42);
\coordinate (v4) at (-0.08,0.12);
\coordinate (v5) at (0.82,0.50);
\coordinate (v6) at (0.97,-0.40);
\coordinate (v7) at (1.88,0.05);
\coordinate (z1) at (-1.80,2.28);
\coordinate (z2) at (0.10,2.28);
% ---------------------------------------------------------------- every z_i is joined to every vertex of G
\foreach \i in {1,...,7}{
  \draw[fan] (z1) -- (v\i);
  \draw[fan] (z2) -- (v\i);}
% ---------------------------------------------------------------- z_1 -- z_2
\draw[zz] (z1) -- (z2);
% ---------------------------------------------------------------- edges and vertices of G
\draw[ge] (v1) -- (v2);
\draw[ge] (v1) -- (v3);
\draw[ge] (v2) -- (v3);
\draw[ge] (v2) -- (v4);
\draw[ge] (v3) -- (v4);
\draw[ge] (v4) -- (v5);
\draw[ge] (v4) -- (v6);
\draw[ge] (v5) -- (v6);
\draw[ge] (v5) -- (v7);
\draw[ge] (v6) -- (v7);
\foreach \i in {1,...,7}{\node[gv] at (v\i) {};}
\node[univ] at (z1) {$z_1$};
\node[univ] at (z2) {$z_2$};
% ---------------------------------------------------------------- labels
\node[font=\small,inner sep=0pt,anchor=base] at (0,-1.50) {$G$:\ \ hard instance for $3$ colours};
\node[note] (lab) at (1.02,2.15) {weight $d_v/3$ to\\[-0.3mm]every vertex $v$};
\draw[ptr] ([xshift=-1.5pt,yshift=-3pt]lab.west) -- (0.50,1.60);
\end{tikzpicture}
\par
\vspace{2.6mm}
The new vertices $z_i$ are\\
\hl{forced to take the extra colours}\par
\vspace{0.8mm}
\paperref{Lemma 6.1}
\end{column}
\begin{column}{0.55\textwidth}
\raggedright
\begin{itemize}\setlength{\itemsep}{1.1mm}
\item \hl{Preserves perfect completeness}
\item Stated constants: applied on top of a weaker intermediate bound
\item Applied directly to $\frac{16}{17}$ \src{AOTW14}:
\end{itemize}
\vspace{0.7mm}
\hspace*{6mm}\pIVbox{$1-\dfrac{3}{51+34L+3\binom L2}$}\hfill
\begin{tabular}{@{}r@{}}{\small\color{soft}$k=3+L$}\\[0.6mm]\paperref{Theorem 1.3}\end{tabular}\par
\vspace{1.5mm}
\hspace*{6mm}{\renewcommand{\arraystretch}{1.1}%
\begin{tabular}{@{}l@{\hspace{4mm}}c@{\hspace{4mm}}l@{}}
\toprule
$4$-colourable & \hl{$82/85$} & {\footnotesize\color{soft}stated before: $85/86$}\\
$5$-colourable & \hl{$119/122$} & {\footnotesize\color{soft}stated before: $120/121$}\\
\bottomrule
\end{tabular}}\par
\vspace{1.4mm}
\hspace*{6mm}{\small\color{soft}Improves the stated bounds for $4\le k\le20$}\par
\vspace{0.4mm}
\hspace*{6mm}{\small\color{soft}Applied to Theorem 1.1: only $47/49$ for $k=4$}\par
\end{column}
\end{columns}
\end{frame}
